% !Mode:: "TeX:UTF-8"
\chapter{\LaTeX{}基础语法}
\label{chapter-basic}
\section{字体字号设置}
\label{sec-font}
\subsection{字体设置}
\label{subsec-font}
本模板使用了CTEX封装的\texttt{ctexbook}类，作为毕设论文，一般不需要
过多使用字体。因此，全文默认字体为~{\bf 宋体}~和{\bf Times New Roman}字体。
默认情况下不需要对字体进行特别的设置，论文中有需要特别强调的地方，可以按表\ref{tab-fonts}
中的相应命令进行设置即可。

\begin{table}[ht]
\begin{center}
\caption{字体的使用}\label{tab-fonts}
\begin{tabular}{c|c|c}
\hline
字体 & 使用命令范例 & 编译效果\\\hline\hline
宋体 & \verb|{\songti|宋体使用范例\} & {\songti 宋体使用范例}\\\hline
黑体 & \verb|{\heiti|黑体使用范例\} & {\heiti 黑体使用范例}\\\hline
加粗字体 & \verb|{\bf|加粗字体Example\} & {\bf 加粗字体Example}\\\hline
加下划线 & \verb|\underline{|加下划线Example\} & \underline{加下划线Example}\\\hline\hline
Roman &\verb|\textrm{|Roman字体范例\} & \textrm{Roman字体范例}\\\hline
Typewriter & \verb|\texttt{|Typewriter字体范例\} & \texttt{Typewriter字体范例}\\\hline
Slanted & \verb|\textsl{|Slanted字体范例\} &\textsl{Slanted字体范例}\\\hline
Bold face & \verb|\textbf{|Bold face字体范例\} & \textbf{Bold face字体范例}\\\hline
Emphasized & \verb|\emph{|Emphasized字体范例\} & \emph{Emphasized字体范例}\\\hline
Document font & \verb|\textnormal{|Document font字体范例\} & \textnormal{Document font字体范例}\\\hline
\end{tabular}
\end{center}
\end{table}

\subsection{字号设置}
\label{subsec-fontsize}
论文要求正文字号为小四号，除非有特别的要求，{\heiti 全文不需要}对字号进行另外的设置。
需要进行字号设置时，可以参照表\ref{tab-fontsize}进行相应的设置。可以从小五号到初号进行设置。

\begin{table}[ht]
\begin{center}
\caption{字号的设置}\label{tab-fontsize}
\begin{tabular}{c|c|c|c}
\hline
使用范例 & 编译效果 &使用范例 & 编译效果\\\hline\hline
\verb|{\zihao{-5}|小五号\} & {\zihao{-5}小五号} & \verb|{\zihao{5}|五号字\} & {\zihao{5}五号字}\\\hline
\verb|{\zihao{-4}|小四号\} & {\zihao{-4}小四号} & \verb|{\zihao{4}|四号字\} & {\zihao{4}四号字}\\\hline
\dots & \dots & \dots & \dots \\\hline
\verb|{\zihao{-0}|小初号\} & {\zihao{-0}小初号} & \verb|{\zihao{0}|初号字\} & {\zihao{0}初号字}\\\hline
\end{tabular}
\end{center}
\end{table}

\subsection{间距间隔控制}
\label{subsec-space}
对于间距的控制，\LaTeX{}中的间距比Office要强大得多了！
Office除了空格回车外，差不多就没有了。而对于\LaTeX{}来说，
可使用的间隔控制命令有如表\ref{tab-space}所示若干，虽然很多都用不到。

\begin{table}
\begin{center}
\caption{\LaTeX{}中的一些常用间隔间距控制命令}\label{tab-space}
\begin{tabular}{l|l}
\hline
命令&效果\slash 说明\\\hline\hline
间隔\texttt{\textbackslash quad}间隔 & 间隔\quad 间隔\\\hline
间隔\texttt{\textbackslash qquad}间隔 & 间隔\qquad 间隔\\\hline
间隔$\sim$间隔 & 间隔~间隔\\\hline
间隔\verb*|\ |间隔 & 间隔\ 间隔\\\hline
间隔\texttt{\textbackslash ,}间隔 & 间隔\, 间隔\\\hline\hline
\texttt{\textbackslash hfill} & 水平填充\\\hline
\texttt{\textbackslash vfill} & 垂直填充\\\hline
\texttt{\textbackslash hspace\{5cm\}} & 水平空白5cm的距离\\\hline
\texttt{\textbackslash vspace\{5cm\}} & 垂直空白5cm的距离\\\hline
\end{tabular}
\end{center}
\end{table}

\LaTeX{}支持的单位有：\texttt{in}---英寸(inch)，\texttt{mm}---毫米，
\texttt{cm}---厘米，\texttt{pt}---Points(大约1\slash 72~inch)，
以及相对距离：\texttt{em}---当前字体的字符“\textbf{M}”的宽度，
\texttt{ex}---当前字体的字符“\textbf{x}”的宽度，长度可以为负值，如~-1.5em~。

换行有两种方式，分别为命令“\verb|\\ |”以及命令\verb|\newline |，
这两个命令将使“光标”切换到下一行，但并不具有缩进（即\LaTeX{}并不认为
这是新的一个段落，当然可以使用\verb|\qquad |来强制缩进两个字符间距）。

而需要另起新的段落时，可以在段落加上\verb|\par |来实现，也可以
通过一个（或多个\footnote{\LaTeX{}会忽略掉.tex文件中的空白距离
和多余的空行。}）空行来达到此目的。具体使用哪种方式，请论文作者自行选择。

\section{章节标题以及条目的设置}
\label{sec-chaptitle}
新的一章将在新的页面开始，使用\texttt{\textbackslash chapter}\{\LaTeX{}基础语法\}
命令时，即可以得到~\hyperref[chapter-basic]{第5章~\LaTeX{}基础语法}~这样的章名称。

同样的，还有命令\textbackslash section\{•\}、\textbackslash subsection\{•\}
来定义节、条的标题。对于章节前面的“第1章”、“1.1”等标号，则不需要用户来编辑，
\LaTeX{}会根据规则来进行编排。

\section{图形}
\label{sec-figure}
在编写论文过程中，需要引用图形甚至多个图形时，可以使用\texttt{figure}环境，
以代码\ref{code-fig}来说明如何使用此环境。

需要注意的是，figure为浮动环境，即\LaTeX{}在排版浮动环境时，
会根据一定的美观规则来排版，并不是将将代码放在某处，图片就会出现在某处，
如代码\ref{code-fig}中的\texttt{[htbp]}为浮动环境选项，
\texttt{h}表示为here，\texttt{t}为top，\texttt{b}为bottom，
\texttt{p}为page，优先顺序依次递减，\LaTeX{}先判断here（当前
位置）是否可以放置此环境（代码\ref{code-fig}中为图片），否则
依次判断top、bottom、page来放置此环境中的内容。

\begin{lstlisting}[label={code-fig},caption={使用figure环境},language={LaTeX}]
 如图 \ref{fig-buaamark}所示为北航校徽。
\begin{figure}[htbp]
\begin{center}
\includegraphics[width=90bp]{figure/buaamark.eps}
\caption{校徽\label{fig-buaamark}}
\end{center}
\end{figure}
\end{lstlisting}

代码中\textbackslash includegraphics后的[width=90bp]为对图片
大小进行控制的代码，可以替换成为height,scale等。
\textbackslash caption为图的标题，\textbackslash label为图的
引用标号，在“如图”后引用即可得到“\ref{figbuaamark}”的效果。
具体生成的效果如第\pageref{figbuaamark}页所示。

如图\ref{figbuaamark}所示为北航校徽。

\begin{figure}[htb]
\begin{center}
\includegraphics[width=90bp]{figure/buaamark.eps}
\caption{校徽\label{figbuaamark}}
\end{center}
\end{figure}

\section{表格}
\label{sec-table}
表格也是论文中经常需要用到的一个元素，它也可以作为一个浮动环境
（当然也可以不浮动）由\LaTeX{}自己来按照一定规则来决定如何放置。

常用的表格浮动环境以代码\ref{code-table}作为示例。

\begin{lstlisting}[caption={表格浮动环境},label={code-table},language={LaTeX}]
\begin{table}[h]
\begin{center}
\caption{浮动表格示例\label{tab:test3}}
\begin{tabular}{c||c|c|c}
\hline
序号&测试内容1&测试内容2&测试内容3\\
\hline \hline
1&一&二&三\\
\hline
2&四&五&六\\
\hline
3&七&八&九\\
\hline
\end{tabular}
\end{center}
\end{table}
\end{lstlisting}

经过\LaTeX{}编译排版出来的结果如表\ref{tab-test3}所示。

\texttt{table}环境即为浮动环境，其后带htbp等参数，其具体
意义在\ref{sec:fig}节描述过了。\texttt{tabular}即为表格环境，
后面的\texttt{\{c||c|c|c\}}的意思为，表格分为四列，第一列
居中{\bf c}\textit{(center)}对齐，其他对齐方式为左
对齐{\bf l}\textit{(left)}，右对齐{\bf r}\textit{(right)}。
然后竖直画两条竖线，第二列和第三列第四列也为居中对齐，并且用一竖线
分隔。表格中需要加横线时使用\texttt{\textbackslash hline}来
添加，横线或竖线可添加多条。

表格中每一项使用“\& ”分隔，每一行使用\texttt{\textbackslash\textbackslash}来进行换行。
\begin{table}
\begin{center}
\caption{浮动表格示例\label{tab-test3}}
\begin{tabular}{c||c|c|c}
\hline
序号&测试内容1&测试内容2&测试内容3\\
\hline \hline
1&一&二&三\\
\hline
2&四&五&六\\
\hline
3&七&八&九\\
\hline
\end{tabular}
\end{center}
\end{table}

\section{数学公式和定理定义等}
\label{sec-math}
数学公式可以说是体现\LaTeX{}强大之处的地方。公式越多，用\LaTeX{}排版出来的论文
效果就越漂亮。

数学公式可分为行间公式和独立\footnote{貌似这个称呼不太准确}公式。{\heiti 行间公式}为放置在
论文段落中间，这种情况一般是将所要编排的公式放置在\verb|\(|和\verb|\)|\footnote{此处
的"("为半角标(英文键盘)时的小括号。}，或者是放置在\verb|$ $|中间，如\verb|$2+2=4$|，
或者是放置在\verb|\begin{math}|和\verb|\end{math}|之间。

而独立公式，即将公式独立于段落放置一行或多行。此种方式又分为编号和非编号公式。
{\heiti 无编号}的独立公式可以使用\verb|\[|和\verb|\]|中间，或者放置在\verb|\begin{displaymath}|
和\verb|\end{displaymath}|中间。而{\heiti 具有编号}的独立公式，
则将公式放置于\verb|\begin{equation}|和\verb|\end{equation}|中间即可。

还有需要注意的一点，\LaTeX{}会认为公式中无空白距离及其他文字出现。如果在公式中需要空白距离时，
可以使用\verb|\quad|、\verb|\qquad|等空白距离命令；需要在公式中使用文本时，
需要将文本放置在\verb|\textrm{...}|中。

其他的一些更深入的数学公式排版知识请查找相应的\LaTeX{}的入门文档或自行搜索。
当然，也有一些工具，可以将繁杂的\LaTeX{}数学代码的书写进行简化。
比较常用的可以使用mathtype将公式导出\LaTeX{}代码。具体请见
附录\ref{chapter-faq}中的问题\ref{mathtype2latex}。

\section{条目及环境}
\label{sec-item}
条目环境分为itemize，enumerate，description三种。itemize用于简单的列表条目，enumerate
用于带序号的列表条目，description用于带描述的列表条目。

以下用代码\ref{code-item}来说明。
\begin{lstlisting}[label={code-item},caption={条目环境},language={LaTeX}]
\begin{enumerate}
    \item 这里是第一条：
        \begin{itemize}
            \item 这里是第一点。
            \item 这里是第二点。
        \end{itemize}
    \item 这里是第二条：
        \begin{description}
            \item[ 观点三 ] 这里是第三点。
            \item[ 观点四 ] 这里是第四点。
        \end{description}
\end{enumerate}
\end{lstlisting}

具体的效果如下所示。
\begin{enumerate}
    \item 这里是第一条：
        \begin{itemize}
            \item 这里是第一点。
            \item 这里是第二点。
        \end{itemize}
    \item 这里是第二条：
        \begin{description}
            \item[观点三] 这里是第三点。
            \item[观点四] 这里是第四点。
        \end{description}
\end{enumerate}

左右对齐时，只需要把需要进行对齐的内容分别放到 flushleft或flushright环境中即可。

\section{使用~listing~显示代码}
\label{sec-listing}
原先~ctexbook~文档类和~listings~宏包配合使用时，代码在换页时会
出现莫名其妙的错误，后来经高人指点，顺利解决了。
感兴趣的话，可以看看\href{http://bbs.ctex.org/viewthread.php?tid=53451}{这里}。
这里给使用~listings~宏包插入源代码的例子，这里是一段C代码。
另外，listings宏包真可谓博大精深，可以实现各种复杂、漂亮的效果，想要进一步学习的同学，可以参考
\href{http://mirror.ctan.org/macros/latex/contrib/listings/listings.pdf}
{listings宏包手册\footnote{此节完全参考武汉理工大学论文模板}}。

\section{更深入一些的操作}
\label{sec-more}
\subsection{字体字号的强调}
\label{subsec-morefont}
可以通过对文字使用\textcolor{blue}{颜}\textcolor{red}{色}来强调说明。
使用\textcolor[rgb]{0.5,0.95,0.96}{有颜色的文字}可以通过
命令\texttt{\textbackslash textcolor}\{red\}\{红色的字\}，
或者\texttt{\textbackslash textcolor[rgb]}\{0.5,0.98,0.7\}\{使用RGB混合颜色\}
来\textcolor[rgb]{0.5,0.98,0.7}{使用RGB混合颜色}进行强调。

\subsection{两个或多个图形的排列}
\label{subsec-morefig}
而对于需要两个或三个图形并排或者矩形排列的时候，则示例代码可以如\ref{code-subfigs}所示。
\begin{lstlisting}[caption={图形的排列},label={code-subfigs},language={LaTeX}]
\begin{figure}[htb]
\begin{center}
    \subfigure[ 并排小图 a]{\label{fig:minia}\includegraphics[width=.2\textwidth]{figure/buaamark.eps}}\hspace{7em}
    \subfigure[ 并排小图 b]{\label{fig:minib}\includegraphics[width=.2\textwidth]{figure/buaamark.eps}}
    \caption{ 多个图形并排的示例 }
    \label{subfig1}
\end{center}
\end{figure}

\begin{figure}[htb]
\begin{minipage}{0.5\textwidth}
    \centering
    \includegraphics[width=0.35\textwidth]{figure/buaamark.eps}
    \caption{ 小图A}
    \label{minifiga}
\end{minipage}
\begin{minipage}{0.5\textwidth}
    \centering
    \includegraphics[width=0.35\textwidth]{figure/buaamark.eps}
    \caption{ 小图B}
    \label{minifigb}
\end{minipage}
\end{figure}

\begin{figure}[htb]
\centering
    \begin{tabular}{cc}
    \subfigure[ 子图 A]{\label{suba}\includegraphics[width=.18\textwidth]{figure/buaamark.eps}}\hspace{7em} & \subfigure[ 子图 B]{\label{subB}\includegraphics[width=.18\textwidth]{figure/buaamark.eps}}\\
    \subfigure[ 子图 C]{\label{subc}\includegraphics[width=.18\textwidth]{figure/buaamark.eps}}\hspace{7em} & \subfigure[ 子图 D]{\label{subd}\includegraphics[width=.18\textwidth]{figure/buaamark.eps}}\\
    \end{tabular}
\caption{ 矩形的 subfig 排列}
\label{maxfig}
\end{figure}
\end{lstlisting}

而具体的效果如下所示。图\ref{subfig1}为子图排列的形式，两个子图有各自的图题，并有一个共同的图题。

\begin{figure}[htb]
\begin{center}
    \subfigure[并排小图a]{\label{fig:minia}\includegraphics[width=.2\textwidth]{figure/buaamark.eps}}\hspace{7em}
    \subfigure[并排小图b]{\label{fig:minib}\includegraphics[width=.2\textwidth]{figure/buaamark.eps}}
    \caption{多个图形并排的示例}
    \label{subfig1}
\end{center}
\end{figure}

图\ref{minifiga}和图\ref{minifigb}为两个并排的图形，各自均有图题。

\begin{figure}[htb]
\begin{minipage}{0.5\textwidth}
    \centering
    \includegraphics[width=0.35\textwidth]{figure/buaamark.eps}
    \caption{小图A}
    \label{minifiga}
\end{minipage}
\begin{minipage}{0.5\textwidth}
    \centering
    \includegraphics[width=0.35\textwidth]{figure/buaamark.eps}
    \caption{小图B}
    \label{minifigb}
\end{minipage}
\end{figure}

图\ref{maxfig}所示为2x2排列的图形，使用\texttt{tabular}的表格环境嵌套在\texttt{figure}环境中来达到此效果。

\begin{figure}[htb]
\centering
    \begin{tabular}{cc}
    \subfigure[子图A]{\label{suba}\includegraphics[width=.18\textwidth]{figure/buaamark.eps}}\hspace{7em} & \subfigure[子图B]{\label{subB}\includegraphics[width=.18\textwidth]{figure/buaamark.eps}}\\
    \subfigure[子图C]{\label{subc}\includegraphics[width=.18\textwidth]{figure/buaamark.eps}}\hspace{7em} & \subfigure[子图D]{\label{subd}\includegraphics[width=.18\textwidth]{figure/buaamark.eps}}\\
    \end{tabular}
\caption{矩形的subfig排列}
\label{maxfig}
\end{figure}
到此，图形的排列基本就OK了。一般只会用到图\ref{figbuaamark}这样的图形
环境就可以了，后面的太复杂，一般情况下不推荐使用。如果对图形的编排的兴趣，
可以google一下《\LaTeX{}插图指南》一书。

\subsection{图形和表格的并排显示}
\label{subsec-morefigtab}
如果需要图形和表格并排显示（一般情况下是不会有这种情况的，也十分不推荐
图形和表格并排显示），则可以通过的代码\ref{code-figtab}来对图形和
表格进行处理来达到并排显示的效果。图\ref{fig-table}和
表\ref{table-fig}就是两个排一起的啦...\par
\begin{figure}[htb]
\begin{minipage}[b]{0.5\textwidth}
\centering
\includegraphics[width=0.6\textwidth]{figure/fgbt.jpg}
\caption{左侧的图片}
\label{fig-table}
\end{minipage}
\begin{minipage}[b]{0.5\textwidth}
\centering
\begin{tabular}{c||c|c} \hline
Day & Data &other\\ \hline\hline
Monday & 1 &1.5\\
Tuesday & 2 &2.5\\
Wednesday & 3 &3.5\\
Thursday & 4 &4.5\\
Friday & 5 &5.5\\ \hline
\end{tabular}
\caption{右侧的表格}
\label{table-fig}
\end{minipage}
\end{figure}
\begin{lstlisting}[label={code-figtab},caption={图形和表格并排显示},language={LaTeX}]
\begin{figure}[htb]
\begin{minipage}[b]{0.5\textwidth}
\centering
\includegraphics[width=0.6\textwidth]{graphics/fgbt.jpg}
\caption{FGBT-left}
\label{fig:by:table}
\end{minipage}%
\begin{minipage}[b]{0.5\textwidth}
\centering
\begin{tabular}{c||c|c} \hline
Day & Data &other\\ \hline\hline
Monday & 1 &1.5\\
Tuesday & 2 &2.5\\
Wednesday & 3 &3.5\\
Thursday & 4 &4.5\\
Friday & 5 &5.5\\ \hline
\end{tabular}
\caption{TABLE-right}
\label{table:by:fig}
\end{minipage}
\end{figure}
\end{lstlisting}

\subsection{长表格的示例}
\label{subsec-moretab}
当需要排版的表格长度超过一页，那么在本模板中推荐使用 \textsf{longtable}
宏包，模板对 \textsf{longtable} 进行了相应的设置，所以用起来可能简单一些。
表~\ref{example-longtable} 就是\textsf{longtable}的简单示例。

\begin{center}
\begin{longtable}{|l|l|l|}
\caption{长表格示例} \label{example-longtable} \\

\hline \multicolumn{1}{|c|}{\textbf{Time (s)}} & \multicolumn{1}{c|}{\textbf{Triple chosen}} & \multicolumn{1}{c|}{\textbf{Other feasible triples}} \\ \hline 
\endfirsthead

\multicolumn{3}{l}%
{{\bfseries \tablename\ \thetable{} --接\,上\,页}} \\
\hline \multicolumn{1}{|c|}{\textbf{Time (s)}} &
\multicolumn{1}{c|}{\textbf{Triple chosen}} &
\multicolumn{1}{c|}{\textbf{Other feasible triples}} \\ \hline
\endhead

\hline \multicolumn{3}{|r|}{{接\,下\,页}} \\ \hline
\endfoot

\hline \hline
\endlastfoot

0 & (1, 11, 13725) & (1, 12, 10980), (1, 13, 8235), (2, 2, 0), (3, 1, 0) \\
2745 & (1, 12, 10980) & (1, 13, 8235), (2, 2, 0), (2, 3, 0), (3, 1, 0) \\
5490 & (1, 12, 13725) & (2, 2, 2745), (2, 3, 0), (3, 1, 0) \\
8235 & (1, 12, 16470) & (1, 13, 13725), (2, 2, 2745), (2, 3, 0), (3, 1, 0) \\
10980 & (1, 12, 16470) & (1, 13, 13725), (2, 2, 2745), (2, 3, 0), (3, 1, 0) \\
13725 & (1, 12, 16470) & (1, 13, 13725), (2, 2, 2745), (2, 3, 0), (3, 1, 0) \\
16470 & (1, 13, 16470) & (2, 2, 2745), (2, 3, 0), (3, 1, 0) \\
19215 & (1, 12, 16470) & (1, 13, 13725), (2, 2, 2745), (2, 3, 0), (3, 1, 0) \\
21960 & (1, 12, 16470) & (1, 13, 13725), (2, 2, 2745), (2, 3, 0), (3, 1, 0) \\
24705 & (1, 12, 16470) & (1, 13, 13725), (2, 2, 2745), (2, 3, 0), (3, 1, 0) \\
27450 & (1, 12, 16470) & (1, 13, 13725), (2, 2, 2745), (2, 3, 0), (3, 1, 0) \\
30195 & (2, 2, 2745) & (2, 3, 0), (3, 1, 0) \\
32940 & (1, 13, 16470) & (2, 2, 2745), (2, 3, 0), (3, 1, 0) \\
35685 & (1, 13, 13725) & (2, 2, 2745), (2, 3, 0), (3, 1, 0) \\
38430 & (1, 13, 10980) & (2, 2, 2745), (2, 3, 0), (3, 1, 0) \\
41175 & (1, 12, 13725) & (1, 13, 10980), (2, 2, 2745), (2, 3, 0), (3, 1, 0) \\
43920 & (1, 13, 10980) & (2, 2, 2745), (2, 3, 0), (3, 1, 0) \\
150975 & (1, 13, 16470) & (2, 2, 2745), (2, 3, 0), (3, 1, 0) \\
153720 & (1, 12, 13725) & (2, 2, 2745), (2, 3, 0), (3, 1, 0) \\
156465 & (1, 13, 13725) & (2, 2, 2745), (2, 3, 0), (3, 1, 0) \\
159210 & (1, 13, 13725) & (2, 2, 2745), (2, 3, 0), (3, 1, 0) \\
161955 & (1, 13, 16470) & (2, 2, 2745), (2, 3, 0), (3, 1, 0) \\
164700 & (1, 13, 13725) & (2, 2, 2745), (2, 3, 0), (3, 1, 0) \\
\end{longtable}
\end{center}

\texttt{longtable}通过\texttt{\textbackslash endfirsthead~~
\textbackslash endhead~~\textbackslash endfoot~~
\textbackslash endlastfoot}来对分页进行相应的控制。
具体可以参照表\ref{example-longtable}及其源代码\ref{code-longtable}
来进行对比分析。

代码中的\texttt{\textbackslash multicolumn}为合并单元格
的命令，具体用法为\texttt{\textbackslash multicolumn \{num$-$cols\}
\{alignment\}\{contents\}}。对于复杂表格中，
\texttt{\textbackslash cline\{i$-$j\}}表示
从第\texttt{i}列开始，画一条横线，直到第\texttt{j}列为止。
具体不清楚用法的话请google吧$\sim$

\begin{lstlisting}[caption={跨页表格示例},label={code-longtable},language={LaTeX}]
\begin{center}
\begin{longtable}{|l|l|l|}
\caption{长表格示例} \label{example:longtable} \\

\hline \multicolumn{1}{|c|}{\textbf{Time (s)}} & \multicolumn{1}{c|}{\textbf{Triple chosen}} & \multicolumn{1}{c|}{\textbf{Other feasible triples}} \\ \hline 
\endfirsthead

\multicolumn{3}{l}%
{{\bfseries \tablename\ \thetable{} --接\,上\,页}} \\
\hline \multicolumn{1}{|c|}{\textbf{Time (s)}} &
\multicolumn{1}{c|}{\textbf{Triple chosen}} &
\multicolumn{1}{c|}{\textbf{Other feasible triples}} \\ \hline
\endhead

\hline \multicolumn{3}{|r|}{{接\,下\,页}} \\ \hline
\endfoot

\hline \hline
\endlastfoot

0 & (1, 11, 13725) & (1, 12, 10980), (1, 13, 8235), (2, 2, 0), (3, 1, 0) \\
2745 & (1, 12, 10980) & (1, 13, 8235), (2, 2, 0), (2, 3, 0), (3, 1, 0) \\
5490 & (1, 12, 13725) & (2, 2, 2745), (2, 3, 0), (3, 1, 0) \\
…… & …… & …… \\
159210 & (1, 13, 13725) & (2, 2, 2745), (2, 3, 0), (3, 1, 0) \\
161955 & (1, 13, 16470) & (2, 2, 2745), (2, 3, 0), (3, 1, 0) \\
164700 & (1, 13, 13725) & (2, 2, 2745), (2, 3, 0), (3, 1, 0) \\
\end{longtable}
\end{center}
\end{lstlisting}

\subsection{定理定义环境}
\label{subsec-morethm}
给大家演示一下各种和证明有关的环境\footnote{本节内容来自武汉理工大学论文模板。}：

\begin{assumption}
待月西厢下，迎风户半开；隔墙花影动，疑是玉人来。
\begin{eqnarray}
  \label{eq:eqnxmp}
  c & = & a^2 - b^2\\
    & = & (a+b)(a-b)
\end{eqnarray}
\end{assumption}

千辛万苦，历尽艰难，得有今日。然相从数千里，未曾哀戚。今将渡江，方图百年欢笑，如
何反起悲伤？（引自《杜十娘怒沉百宝箱》）

\begin{definition}
子曰：「道千乘之国，敬事而信，节用而爱人，使民以时。」
\end{definition}

千古第一定义！问世间、情为何物，只教生死相许？天南地北双飞客，老翅几回寒暑。欢乐趣，离别苦，就中更有痴儿女。
君应有语，渺万里层云，千山暮雪，只影向谁去？

横汾路，寂寞当年箫鼓，荒烟依旧平楚。招魂楚些何嗟及，山鬼暗谛风雨。天也妒，未信与，莺儿燕子俱黄土。
千秋万古，为留待骚人，狂歌痛饮，来访雁丘处。

\begin{proposition}
 曾子曰：「吾日三省吾身 --- 为人谋而不忠乎？与朋友交而不信乎？传不习乎？」
\end{proposition}

多么凄美的命题啊！其日牛马嘶，新妇入青庐，奄奄黄昏后，寂寂人定初，我命绝今日，
魂去尸长留，揽裙脱丝履，举身赴清池，府吏闻此事，心知长别离，徘徊庭树下，自挂东南
枝。

\begin{remark}
天不言自高，水不言自流。
\begin{gather*}
\begin{split} 
\varphi(x,z)
&=z-\gamma_{10}x-\gamma_{mn}x^mz^n\\
&=z-Mr^{-1}x-Mr^{-(m+n)}x^mz^n
\end{split}\\[6pt]
\begin{align} \zeta^0&=(\xi^0)^2,\\
\zeta^1 &=\xi^0\xi^1,\\
\zeta^2 &=(\xi^1)^2,
\end{align}
\end{gather*}
\end{remark}

天尊地卑，乾坤定矣。卑高以陈，贵贱位矣。 动静有常，刚柔断矣。方以类聚，物以群分，
吉凶生矣。在天成象，在地成形，变化见矣。鼓之以雷霆，润之以风雨，日月运行，一寒一
暑，乾道成男，坤道成女。乾知大始，坤作成物。乾以易知，坤以简能。易则易知，简则易
从。易知则有亲，易从则有功。有亲则可久，有功则可大。可久则贤人之德，可大则贤人之
业。易简，而天下矣之理矣；天下之理得，而成位乎其中矣。

\begin{axiom}
两点间直线段距离最短。  
\begin{align}
x&\equiv y+1\pmod{m^2}\\
x&\equiv y+1\mod{m^2}\\
x&\equiv y+1\pod{m^2}
\end{align}
\end{axiom}

《彖曰》：大哉乾元，万物资始，乃统天。云行雨施，品物流形。大明始终，六位时成，时
乘六龙以御天。乾道变化，各正性命，保合大和，乃利贞。首出庶物，万国咸宁。

《象曰》：天行健，君子以自强不息。潜龙勿用，阳在下也。见龙再田，德施普也。终日乾
乾，反复道也。或跃在渊，进无咎也。飞龙在天，大人造也。亢龙有悔，盈不可久也。用九，
天德不可为首也。 　　

\begin{lemma}
《猫和老鼠》是我最爱看的动画片。
\begin{multline*}%\tag*{[a]} % 这个不出现在索引中
\int_a^b\biggl\{\int_a^b[f(x)^2g(y)^2+f(y)^2g(x)^2]
 -2f(x)g(x)f(y)g(y)\,dx\biggr\}\,dy \\
 =\int_a^b\biggl\{g(y)^2\int_a^bf^2+f(y)^2
  \int_a^b g^2-2f(y)g(y)\int_a^b fg\biggr\}\,dy
\end{multline*}
\end{lemma}

行行重行行，与君生别离。相去万余里，各在天一涯。道路阻且长，会面安可知。胡马依北
风，越鸟巢南枝。相去日已远，衣带日已缓。浮云蔽白日，游子不顾返。思君令人老，岁月
忽已晚。  弃捐勿复道，努力加餐饭。

\begin{theorem}\label{the:theorem1}
犯我强汉者，虽远必诛\hfill --- 陈汤（汉）
\end{theorem}
\begin{subequations}
\begin{align}
y & = 1 \\
y & = 0
\end{align}
\end{subequations}
道可道，非常道。名可名，非常名。无名天地之始；有名万物之母。故常无，欲以观其妙；
常有，欲以观其徼。此两者，同出而异名，同谓之玄。玄之又玄，众妙之门。上善若水。水
善利万物而不争，处众人之所恶，故几于道。曲则全，枉则直，洼则盈，敝则新，少则多，
多则惑。人法地，地法天，天法道，道法自然。知人者智，自知者明。胜人者有力，自胜
者强。知足者富。强行者有志。不失其所者久。死而不亡者寿。

\begin{proof}
燕赵古称多感慨悲歌之士。董生举进士，连不得志于有司，怀抱利器，郁郁适兹土，吾
知其必有合也。董生勉乎哉？

夫以子之不遇时，苟慕义强仁者，皆爱惜焉，矧燕、赵之士出乎其性者哉！然吾尝闻
风俗与化移易，吾恶知其今不异于古所云邪？聊以吾子之行卜之也。董生勉乎哉？

吾因子有所感矣。为我吊望诸君之墓，而观于其市，复有昔时屠狗者乎？为我谢
曰：“明天子在上，可以出而仕矣！” \hfill --- 韩愈《送董邵南序》
\end{proof}

\begin{corollary}
  四川话配音的《猫和老鼠》是世界上最好看最好听最有趣的动画片。
\begin{alignat}{3}
V_i & =v_i - q_i v_j, & \qquad X_i & = x_i - q_i x_j,
 & \qquad U_i & = u_i,
 \qquad \text{for $i\ne j$;}\label{eq:B}\\
V_j & = v_j, & \qquad X_j & = x_j,
  & \qquad U_j & u_j + \sum_{i\ne j} q_i u_i.
\end{alignat}
\end{corollary}

迢迢牵牛星，皎皎河汉女。
纤纤擢素手，札札弄机杼。
终日不成章，泣涕零如雨。
河汉清且浅，相去复几许。
盈盈一水间，脉脉不得语。

\begin{example}
  大家来看这个例子。
\begin{equation}
\label{ktc}
\left\{\begin{array}{l}
\nabla f({\mbox{\boldmath $x$}}^*)-\sum\limits_{j=1}^p\lambda_j\nabla g_j({\mbox{\boldmath $x$}}^*)=0\\[0.3cm]
\lambda_jg_j({\mbox{\boldmath $x$}}^*)=0,\quad j=1,2,\cdots,p\\[0.2cm]
\lambda_j\ge 0,\quad j=1,2,\cdots,p.
\end{array}\right.
\end{equation}
\end{example}

\begin{exercise}
  清列出 Andrew S. Tanenbaum 和 W. Richard Stevens 的所有著作。
\end{exercise}

\begin{conjecture} \textit{Poincare Conjecture} If in a closed three-dimensional
  space, any closed curves can shrink to a point continuously, this space can be
  deformed to a sphere.
\end{conjecture}

\begin{problem}
 回答还是不回答，是个问题。 
\end{problem}

如何引用定理~\ref{the:theorem1} 呢？加上 \verb|label| 使用 \verb|ref| 即可。妾发
初覆额，折花门前剧。郎骑竹马来，绕床弄青梅。同居长干里，两小无嫌猜。 十四为君妇，
羞颜未尝开。低头向暗壁，千唤不一回。十五始展眉，愿同尘与灰。常存抱柱信，岂上望夫
台。 十六君远行，瞿塘滟滪堆。五月不可触，猿声天上哀。门前迟行迹，一一生绿苔。苔深
不能扫，落叶秋风早。八月蝴蝶来，双飞西园草。感此伤妾心，坐愁红颜老。